\bookchapter{Глава 25: Утилизация б/у телефона и цифровое наследие}

\chapterintro{Привет, друзья! После двадцати с лишним глав «культивации» вы, верю, выросли из «цифрового новичка» в довольно уверенного пользователя телефона. Но технологическая жизнь циклична — новая любовь рано или поздно становится старой. Когда вы готовитесь «уволить» старого товарища, знаете ли вы, как достойно и безопасно попрощаться?}

\booksection{«Эвтаназия» старого телефона: руководство по безопасной продаже или передаче}

Ключевой принцип: очистка данных — это не «удаление», а «перезапись». Многие думают, что достаточно перетащить фото в корзину и очистить альбом — и дело с концом. Но после перепродажи телефона новый владелец запускает программу восстановления данных — и старые фото, переписки, фото удостоверения все возвращаются. Такой «цифровой позор» случается каждый год.

\begin{enumerate}
  \item \textbf{Генеральная «уборка»}: выполните последний полный бэкап, выйдите из всех аккаунтов (iCloud, Huawei/Xiaomi/Samsung, Google), отключите «Найти мой телефон» — иначе следующий пользователь будет заблокирован активационным замком.
  \item \textbf{Финальная «форматизация»}: «Восстановление заводских настроек» на современных телефонах обычно включает криптографическое стирание — уничтожается сам ключ шифрования; без ключа все данные мгновенно становятся нечитаемыми.
  \item \textbf{Последний «обманный манёвр» (только для старых устройств)}: современные телефоны (iOS 8+, Android 6+) по умолчанию используют полное дисковое шифрование (FBE); восстановление заводских настроек уничтожает ключ шифрования, и данные теоретически невосстановимы —\textbf{этот шаг не нужен}. Только для очень старых нешифрованных устройств (до 2015 года) после сброса нужно заполнить хранилище «мусорными» файлами и снова выполнить сброс.
\end{enumerate}

\textbf{Справочная таблица путей стирания данных по брендам}: принцип один и тот же, но каждый бренд прячет «восстановление заводских настроек» в разных местах; ключевые отличия:

\begin{center}
\begin{tabularx}{\linewidth}{|X|X|X|}
\hline
\textbf{Бренд} & \textbf{Путь стирания} & \textbf{Ключевые отличия и моменты, требующие внимания} \tabularnewline
\hline
Huawei & Настройки Система и обновление Сброс Восстановление заводских настроек & Нужно сначала выйти из аккаунта Huawei и отключить «Найти мой телефон», иначе сработает активационный замок; в режиме PureMode может потребоваться ввод пароля. \tabularnewline
Xiaomi / Redmi & Настройки Моё устройство Восстановление заводских настроек & Внизу опции «Форматировать виртуальную SD-карту» и «Очистить все данные и настройки» — обе нужно отметить; перед выходом из аккаунта Xiaomi отключите «Найти устройство». \tabularnewline
OPPO / OnePlus & Настройки Другие настройки Сброс телефона Полностью очистить все данные & ColorOS предлагает режим «полная очистка» — многократно перезаписывает хранилище; требует ввода пароля блокировки экрана. \tabularnewline
vivo & Настройки Управление системой Бэкап и сброс Очистить все данные & Нужно отметить «Форматировать хранилище телефона», иначе SD-область не очистится; нужно выйти из аккаунта vivo и отключить «Найти телефон». \tabularnewline
Samsung & Настройки Общее управление Сброс Восстановление заводских настроек & Нужно сначала удалить аккаунт Samsung и отключить «Найти мой телефон»; на некоторых моделях есть опция «автоперезапись» — рекомендуется включить. \tabularnewline
iPhone & Настройки Основные Перенос или сброс iPhone Стереть все содержимое и настройки & Обязательно сначала выйдите из Apple ID, иначе следующий пользователь будет заблокирован iCloud активационным замком; iOS 15+ поддерживает «сохранить данные iCloud при выходе» для последующей проверки. \tabularnewline
\hline
\end{tabularx}
\end{center}

\begin{tcbwarn}{Три вещи, которые чаще всего забывают}
\begin{itemize}
  \item \textbf{Не вышел из аккаунта}:после продажи телефон удалённо блокируется, и вы сами не можете войти — в итоге приходится идти с паспортом в официальный сервис для разблокировки.
  \item \textbf{Не извлёк SD-карту}:сброс к заводским настройкам не форматирует фото на внешней SD-карте, не забудьте извлечь её вручную или отформатировать отдельно.
  \item \textbf{Не отключил USB-накопитель / OTG-периферию}:если в процессе стирания подключены внешние устройства, можно случайно удалить данные на них, или процесс стирания прервётся.
\end{itemize}

\end{tcbwarn}

\booksection{II. Цифровое наследство: как сохранить вашу социальную душу?}

Данные в телефоне давно перестали быть «виртуальными». Они хранят ваши деньги, социальные связи, результаты творчества и даже вашу «цифровую личность».\textbf{Что происходит с этими активами после смерти человека?}Это больше не абстрактный вопрос — всё больше семей сталкиваются с этой реальной проблемой: родственники не знают пароли, платформы не разрешают наследование, криптовалюта теряется навсегда.

\begin{tcbinfo}{Правовое основание: статья 49 Закона о защите персональной информации}
Статья 49 Закона КНР «О защите персональной информации», вступившего в силу в 2021 году, прямо устанавливает:\textbf{в случае смерти физического лица его близкие родственники вправе, в целях своих законных и правомерных интересов, осуществлять предусмотренные настоящей главой права в отношении персональной информации умершего — право на ознакомление, копирование, исправление и удаление; за исключением случаев, когда умерший при жизни распорядился иначе.}

Это означает:

\begin{itemize}
  \item \textbf{Близкие родственники имеют правовое основание}для обращения к платформе с заявлением о обработке аккаунта и персональной информации умершего.
  \item «Законные и правомерные интересы» обычно включают оформление наследства, закрытие аккаунта, удаление конфиденциальной информации и т.д.
  \item Если умерший при жизни оставил завещание или договор с иными распоряжениями (например, указал контактное лицо для наследства), выполняются его указания.
  \item Платформа не может под предлогом «конфиденциальности пользователя» полностью отказать в сотрудничестве, однако конкретные процедуры у каждой платформы свои.
\end{itemize}

\end{tcbinfo}

\booksubsection{WeChat и Alipay: аккаунт нельзя унаследовать, но деньги — можно}

У этих двух самых популярных в Китае приложений правила разные:

\begin{itemize}
  \item \textbf{Аккаунт WeChat}:согласно «Лицензионному соглашению на программное обеспечение и услуги Tencent WeChat», право собственности на аккаунт принадлежит Tencent, пользователь имеет лишь право пользования,\textbf{сам аккаунт нельзя унаследовать или передать}. Однако средства на аккаунте —\textbf{баланс и деньги в «Lingqiantong» (WeChat Pay)}— являются личным имуществом; близкие родственники могут предоставить свидетельство о смерти, подтверждение родства и нотариальное свидетельство о праве на наследство для вывода или перевода средств через Tenpay.
  \item \textbf{Alipay}:аккаунт также нельзя напрямую унаследовать, но баланс Alipay, Yu'ebao, Li Caitong и другие\textbf{средства являются наследуемым имуществом}. Наследники могут обратиться в «горячую линию поддержки» Alipay или лично подать свидетельство о смерти и нотариальное свидетельство о наследстве; финансовые продукты Ant Wealth требуют прохождения процедур соответствующей управляющей компании фонда.
  \item \textbf{История переписки}:юридически является личной тайной, платформа не предоставляет её семье по своей инициативе. Если переписка содержит ключи к наследству, можно запросить данные через судебный иск или нотариуса.
\end{itemize}

\begin{tcbwarn}{Чего категорически нельзя делать}
\begin{itemize}
  \item \textbf{Не передавайте пароль от аккаунта умершего семье для «продолжения использования»}— это нарушает соглашение с платформой; при блокировке средств или срабатывании контроля риска доказать правоту будет гораздо сложнее.
  \item \textbf{Не пытайтесь получить код верификации на телефон умершего для «вывода средств»}— это может быть расценено как незаконное присвоение, что уже является правонарушением.
  \item \textbf{Не откладывайте}:некоторые аккаунты при длительном отсутствии входа автоматически удаляются, и возврат средств станет ещё сложнее.
\end{itemize}

\end{tcbwarn}

\booksubsection{«Контактное лицо для наследства» Apple: заранее укажите — и получите вдвое больше результата}

Начиная с iOS 15.2 Apple представила функцию «Legacy Contact», позволяющую заранее указать доверенное контактное лицо; после кончины это лицо может по\textbf{свидетельству о смерти + ключу доступа контактного лица для наследства}запросить доступ к данным iCloud умершего (фото, заметки, контакты и т.д.). Порядок действий:

\begin{enumerate}
  \item На iPhone перейдите в\textbf{[Настройки] > [Ваше имя] > [Вход и безопасность] > [Контактное лицо для наследства]}.
  \item Нажмите «Добавить контактное лицо для наследства», подтвердите личность с помощью Face ID.
  \item Выберите контакт (из адресной книги), система сгенерирует\textbf{ключ доступа}.
  \item Передайте ключ доступа через SMS, электронную почту или распечатайте на бумаге и передайте лично, напомнив о необходимости надёжного хранения.
  \item Можно добавить несколько контактных лиц как резервных, чтобы избежать единой точки отказа.
\end{enumerate}

Указанному лицу не нужно устройство Apple — достаточно ключа и свидетельства о смерти, чтобы подать заявку на доступ на сайте Apple. Весь процесс обычно занимает до 7 дней.

\booksubsection{«Менеджер неактивных аккаунтов» Google: автоматическая обработка дел после вас}

Google предоставляет «Inactive Account Manager», который может автоматически выполнить заданные действия, если пользователь долго не входил в аккаунт — это одна из самых продуманных систем цифрового наследования на сегодняшний день.

\begin{enumerate}
  \item Перейдите на\textbf{myaccount.google.com/inactive}, войдите в аккаунт Google.
  \item Установите тайм-аут: срабатывает через 3 / 6 / 12 / 18 месяцев без входа.
  \item Добавьте доверенных контактных лиц (до 10 человек) и укажите, к каким данным каждый получит доступ (Gmail, Фото, Диск и т.д.).
  \item Опционально: настройте автоответчик, чтобы сообщить отправителям, что вы больше не можете отвечать.
  \item Наконец, выберите: по истечении срока —\textbf{поделиться данными с контактами}или\textbf{удалить аккаунт}.
  \item Укажите номер телефона для дополнительной верификации, чтобы избежать ложного срабатывания.
\end{enumerate}

\booksubsection{Аккаунты в китайских соцсетях: мемориализация и удаление}

Крупнейшие китайские социальные платформы подходят к аккаунтам умерших двумя путями:\textbf{«мемориализация» (сохранение аккаунта для памяти родных и друзей)}и\textbf{«удаление» (полная очистка)}. Родственники могут подать заявление в службу поддержки платформы, предоставив свидетельство о смерти, подтверждение родства и удостоверение личности заявителя.

\begin{itemize}
  \item \textbf{Weibo}:поддерживает «мемориализацию» аккаунта — родственники подают свидетельство о смерти, аккаунт блокируется для входа, но содержимое сохраняется; на главной странице появляется пометка «аккаунт переведён в мемориальный режим». Также можно подать на удаление, требуется нотариальное свидетельство о наследстве.
  \item \textbf{Douyin}:через «Профиль - Настройки - Обратная связь и помощь - Аккаунт и безопасность - Удаление аккаунта» или через службу поддержки вручную, предоставив свидетельство о смерти, подтверждение родства и удостоверение личности заявителя; после проверки аккаунт можно удалить или перевести в мемориальный статус.
  \item \textbf{Bilibili}:через службу поддержки подать доказательства кончины и родства; можно запросить удаление аккаунта или сохранение как мемориальной страницы; обработка загруженных UP-автором видео зависит от авторских прав.
  \item \textbf{Xiaohongshu}:через официальную поддержку подать апелляцию с свидетельством о смерти и подтверждением родства; можно запросить удаление аккаунта и контента.
\end{itemize}

\begin{tcbinfo}{Ключевые напоминания}
\begin{itemize}
  \item Правила платформ обновляются — перед оформлением ознакомьтесь с актуальными инструкциями службы поддержки.
  \item Все подтверждающие документы рекомендуется сначала отсканировать и сохранить копии, а оригиналы хранить надёжно.
  \item При массовом оформлении на нескольких платформах можно один раз получить\textbf{нотариальное свидетельство о праве на наследство}, чтобы не ходить по инстанциям многократно.
\end{itemize}

\end{tcbinfo}

\booksubsection{Проблема наследования цифровых активов}

Помимо социальных аккаунтов, всё больше людей владеют\textbf{действительно ценными цифровыми активами}, наследование которых ещё сложнее:

\begin{itemize}
  \item \textbf{Криптовалюта (Bitcoin, Ethereum и др.)}:по принципу «приватный ключ = право собственности»\textbf{без приватного ключа активы считаются утраченными}. Наследники не могут «подать апелляцию» в блокчейн через юридические документы. Рекомендация: разделите мнемоническую фразу / приватный ключ на фрагменты и храните в нескольких физически безопасных местах (например, банковская ячейка + домашний сейф), а в завещании укажите способ доступа.
  \item \textbf{Аккаунт на бирже}:аналогично банковскому счёту, можно подать нотариальное свидетельство о наследстве бирже (например, Binance, OKX) для распоряжения средствами, однако у зарубежных бирж процесс длительный.
  \item \textbf{Доменное имя}:является наследуемым имуществом; наследники могут подать нотариальное свидетельство регистратору (например, Aliyun, GoDaddy) для переоформления.
  \item \textbf{Интеллектуальная собственность (оригинальные статьи, видео, патенты, авторские права на ПО)}:имущественные авторские права наследуются (личные неимущественные — нет), право на доход переходит к наследникам на 50 лет после смерти автора.
  \item \textbf{Облачные хранилища и сетевые диски}:данные контента принадлежат лично и наследуемы; однако некоторые платформы в соглашении запрещают передачу аккаунта, поэтому нужно сначала скачать данные, затем удалить аккаунт.
\end{itemize}

\textbf{Реальная цена «исчезновения» криптовалюты}\par

По данным Chainalysis, в мире\textbf{более 3 миллионов биткоинов}навсегда потеряны из-за утраты приватных ключей — это примерно одна шестая часть от текущего общего объёма. Если вы владеете криптоактивами, убедитесь, что хотя бы один человек, которому вы абсолютно доверяете, знает\textbf{путь восстановления}— иначе эти богатства навсегда останутся спящими на блокчейне.

\booksubsection{Шаблон «цифрового завещания»}

Традиционное завещание не охватывает цифровые активы. Рекомендуется ежегодно обновлять приложение «цифровое завещание» и хранить его вместе с официальным завещанием. Ниже — шаблон для справки:

\textbf{«Приложение к распоряжению цифровыми активами» (шаблон)}\par

\textbf{Часть I: Список аккаунтов}(записывайте только имя аккаунта и платформу,\textbf{без паролей})

\begin{itemize}
  \item Социальные: WeChat (аккаунт xxx), Weibo, Douyin, Bilibili, Xiaohongshu и др.
  \item Платёжные: Alipay, WeChat Pay, приложения банковских карт.
  \item Облачные: iCloud, Huawei Cloud, Mi Cloud, Google, Microsoft.
  \item Электронная почта: QQ Mail, NetEase Mail, Gmail и др.
  \item Прочие: сетевые диски, подписки, игровые аккаунты и т.д.
\end{itemize}

\textbf{Часть II: Схема хранения паролей}

\begin{itemize}
  \item Способ хранения мастер-пароля: использовать менеджер паролей (например, Bitwarden, 1Password), мастер-пароль записать на бумаге и хранить в домашнем сейфе.
  \item Ключи восстановления / мнемонические фразы: хранить фрагментированно, минимум в двух физических местах (например, банковская ячейка + домашний сейф).
  \item Уполномоченный для запросов: ФИО, номер удостоверения личности, степень родства, контактные данные.
\end{itemize}

\textbf{Часть III: Воля по распоряжению}

\begin{itemize}
  \item Для каждого типа аккаунта — «мемориализация» или «удаление».
  \item Указаны ли для данных iCloud / Google «контактное лицо для наследства» или «менеджер неактивных аккаунтов».
  \item Кто наследует криптоактивы; описание пути доступа к приватному ключу.
  \item Кому принадлежит право на доход от авторских работ.
\end{itemize}

\textbf{Часть IV: Контактные лица и платформы}

\begin{itemize}
  \item FIO доверенного исполнителя, контакты, резервные лица.
  \item Входы в службы поддержки платформ и список требуемых подтверждающих документов.
  \item Дата обновления этого приложения: \_\_\_\_г.\_\_\_м.\_\_\_д., следующая проверка: \_\_\_\_г.\_\_\_м.\_\_\_д.
\end{itemize}

\begin{tcbwarn}{Юридическое замечание о «паролях»}
Вписывать пароли прямо в завещание рискованно — после вскрытия завещания все наследники их увидят.\textbf{Надёжнее: в завещании указывайте только «список аккаунтов + волю по распоряжению + место хранения паролей», а сами пароли передавайте доверенному исполнителю отдельно через мастер-пароль менеджера паролей}. Это соответствует юридическим требованиям и снижает риск утечки паролей.

\end{tcbwarn}

\booksection{III. Экологичный финал}

В одном телефоне содержатся десятки видов металлов — золото, серебро, медь, редкоземельные элементы и др.; выбрасывание на свалку загрязняет почву и расходует ресурсы. Правильных каналов утилизации три, каждый со своими особенностями:

\begin{center}
\begin{tabularx}{\linewidth}{|X|X|X|X|}
\hline
\textbf{Канал} & \textbf{Цена выкупа} & \textbf{Безопасность данных} & \textbf{Подходящий сценарий} \tabularnewline
\hline
Официальный Trade-in (Huawei/Apple/Xiaomi и др.) & Средняя, часто с бонусными купонами & Наивысшая, стандартизированный процесс & Выгоднее всего при покупке нового устройства того же бренда \tabularnewline
Сторонние платформы выкупа (Aihuishou, Zhuanzhuan, Huishoubao и др.) & Несколько выше, можно сравнивать & Высокая, есть отчёт об очистке данных & Хотите сразу получить деньги, без привязки к бренду \tabularnewline
Благотворительная передача (сельские школы, НПО) & Без финансовой отдачи & Нужно полностью стереть самостоятельно & Устройство ещё работоспособно, хотите сделать доброе дело \tabularnewline
Вторичная жизнь (переделка для себя) & Экономия на новом устройстве & Владелец продолжает использовать, без риска & Устройство ещё в норме, дома есть потребность в переделке \tabularnewline
\hline
\end{tabularx}
\end{center}

\begin{itemize}
  \item \textbf{Официальный выкуп / Trade-in}:самый рекомендуемый способ, стандартизированный процесс, безопасность данных гарантирована. На сайтах Apple, Huawei, Xiaomi и др. есть вход в программу Trade-in; после оценки сумма засчитывается в стоимость нового устройства.
  \item \textbf{Лицензированная платформа выкупа}:выбирайте крупную платформу с хорошей репутацией; после заказа курьер SF Express заберёт устройство, после проверки средства переводятся. Обязательно выбирайте платформу,\textbf{предоставляющую «отчёт об очистке данных»}, и сохраняйте подтверждение.
  \item \textbf{Благотворительная передача}:свяжитесь с местной сельской школой или общественной организацией; старый телефон можно переделать в учебный планшет или телефон для пожилых.\textbf{Перед передачей обязательно выполните полный цикл стирания данных и извлеките все SIM/SD-карты}.
  \item \textbf{Вторичная жизнь}:переделайте в автомобильный навигатор, центр умного дома, камеру наблюдения, радионяню, электронную фоторамку. Старый телефон может прослужить ещё несколько лет.
\end{itemize}

\begin{tcbinfo}{Правда об электронных отходах}
Печатная плата телефона содержит\textbf{золото, серебро, палладий, медь}и другие драгоценные металлы; содержание золота в тонне отработанных плат в десятки раз выше, чем в тонне золотой руды. Но при выбрасывании на свалку свинец, ртуть и кадмий просачиваются в почву и грунтовые воды, загрязняя их на десятилетия.\textbf{Правильная утилизация защищает окружающую среду и позволяет дефицитным металлам возвращаться в оборот.}

\end{tcbinfo}

\begin{tcbpractice}{Практика: устроьте вашему старому телефону «стандартные похороны»}
\begin{enumerate}
  \item Последняя полная резервная копия: загрузите фото в облако, скопируйте важные файлы на компьютер, экспортируйте историю чатов.
  \item Выйдите из всех аккаунтов: iCloud / Huawei / Xiaomi / Google, отключите «Найти мой телефон».
  \item Извлеките SIM-карту и SD-карту, отключите всю OTG-периферию.
  \item Выполните шифрованный сброс к заводским настройкам; при необходимости вручную отметьте «форматировать хранилище».
  \item (Только для старых устройств) Заполните хранилище мусорными файлами, затем снова выполните сброс к заводским настройкам.
  \item Выберите следующий этап: официальный Trade-in, лицензированная платформа выкупа, благотворительная передача или передача семье.
  \item Обновите «реестр цифрового наследства», вычеркнув списанные устройства.
\end{enumerate}

\end{tcbpractice}

\begin{tcbpractice}{Список задач на этот месяц}
\begin{enumerate}
  \item Для вашего основного телефона укажите «контактное лицо для наследства» Apple или настройте «менеджер неактивных аккаунтов» Google.
  \item С помощью шаблона из этого раздела составьте «Приложение к распоряжению цифровыми активами» и обсудите с семьёй место хранения.
  \item Проверьте, есть ли дома 2+ простаивающих старых телефона, и запишитесь на выкуп через официальный или лицензированный канал.
\end{enumerate}

\end{tcbpractice}
